\begin{idstatement}{C.2.INT.1}{Proved elliptic-image restriction}
\label{st:C.2.INT.1}
Let $Y$ be a smooth compact complex threefold with
$H^2(Y,\Q)=0$, let $S_i$ be an irreducible component of
a connected reducible divisor $S$, and let
$h:Z\to S_i\hookrightarrow Y$ be the minimal resolution
of its normalization. Suppose that the minimal model
$\sigma:Z\to Z_0$ is a Hopf surface or a non-K\"ahler
elliptic surface. Every strict transform $E\subset Z$ of
an elliptic curve on $Z_0$ has a one-dimensional image
under $h$, and the normalization of that image is elliptic.
The actual restriction homomorphism
$\Pic^0(Y)\to\Pic^0(E)$ is nonconstant for each such $E$.
In particular such a component cannot meet a component
whose minimal model is an Enoki surface.
\end{idstatement}
\begin{proof}
Choose an adjacent component $S_j$. The actual pullback
of its canonical divisor section gives a nonzero effective
Cartier divisor
\[
 A=h^*S_j>0,\qquad \OO_Z(A)=h^*\OO_Y(S_j).
 \tag{C.INT.1}
\]
Its first Chern class is torsion. For some $q>0$,
$\OO_Y(qS_j)\in\Pic^0(Y)$, so the restriction homomorphism
$\rho:\Pic^0(Y)\to\Pic^0(Z)$ has in its image the line
$\OO_Z(qA)$. That line has strictly positive degree for
every Gauduchon metric on $Z$.

We spell out the blowdown comparison. Put $A_0=\sigma_*A$.
The difference $A-\sigma^*A_0$ is exceptional and has
zero real Chern class: both $A$ and its pushforward have
zero real class. The negative definite exceptional
intersection matrix makes every coefficient of this
difference zero. Thus
\[
 A=\sigma^*A_0,\qquad A_0>0.
 \tag{C.INT.2}
\]
Nonzeroness follows by the same negative definiteness if
$A_0$ were zero. Pullback is an isomorphism
$\Pic^0(Z_0)\simeq\Pic^0(Z)$; the integral blowup
formula and $R\sigma_*\OO_Z=\OO_{Z_0}$ identify both
the lattice and the coefficient space defining it.
We therefore regard $\rho$ as a map into $\Pic^0(Z_0)$.
Its image contains $\OO_{Z_0}(qA_0)$.

We first prove the stronger assertion that the composite
restriction to $E$ cannot be zero. Suppose that it is zero.
This situation would in particular arise if $h(E)$ were
a point or had rational normalization.
In the latter case the map from the normal curve $E$
factors as $E\to\PP^1\to Y$. In both cases every line
from $\Pic^0(Y)$ restricts trivially to $E$. Here the
pullback to $\PP^1$ is in $\Pic^0(\PP^1)=0$; this is
the specified map giving the vanishing. The restriction
of $\sigma$ to $E$ is an isomorphism onto its smooth
elliptic image $E_0$.

For a Hopf $Z_0$ the homomorphism
\[
 H^1(Z_0,\R)\longrightarrow H^1(E_0,\R)
 \tag{C.INT.3}
\]
is nonzero. On a primary Hopf surface this follows from
its universal cover $\C^2\setminus\{0\}$: a finite image
of $\pi_1(E_0)$ in $\pi_1(Z_0)=\Z$ would lift a finite
unramified cover of $E_0$ to that cover. Its two coordinate
functions would be constant, contradicting the nonconstant
inclusion of the curve. For a secondary Hopf surface,
pull back along a finite unramified primary cover. A
component above $E_0$ is an elliptic curve. Transfer makes
$H^1(Z_0,\R)\to H^1(Z'_0,\R)$ an isomorphism between
one-dimensional spaces, and the preceding primary
argument proves (C.INT.3) after this pullback, hence
before it. The real-to-holomorphic coefficient injection
of D.1.7, applied on $E_0$, shows that
$H^1(Z_0,\OO)\to H^1(E_0,\OO)$ is nonzero. Since
$\Pic^0(Z_0)=\C^*$, the identity component of the
kernel of $\Pic^0(Z_0)\to\Pic^0(E_0)$ is zero.
Connectedness of $\Pic^0(Y)$ would therefore make
$\rho$ zero, contradicting (C.INT.1)--(C.INT.2).

For a non-K\"ahler elliptic $p:Z_0\to B$, every curve
on its minimal model is a reduced elliptic fibre.
The actual covering construction of
\cite[Lemma~2.2 and \S2]{BrinzanescuMoraru2005}
provides a finite cyclic map $\psi:Z'_0\to Z_0$ and a
principal elliptic bundle $p':Z'_0\to B'$ such that a
reduced fibre $E'$ above $E_0$ maps to $E_0$ by a finite
isogeny. To include multiple fibres, take a base cover
whose branching orders are divisible by their
multiplicities and normalize the fibre product. In the
local cyclic presentation of a smooth multiple fibre,
this removes the transverse multiplicity and leaves a
translation quotient on the reduced fibre, giving exactly
the stated isogeny. Extra branch points over ordinary
fibres make the total map branched there, but its map on
the reduced elliptic fibre is an isomorphism.

On this principal bundle the source's explicit
presentation is
\[
 Z'_0=\Theta^*/\langle\tau\rangle,
 \quad |\tau|>1,\quad \deg\Theta>0,
 \qquad
 \Pic^0(Z'_0)=p'^*\Pic^0(B')\times\C^*.
 \tag{C.INT.4}
\]
In the last factor a scalar $\alpha$ is the automorphy
factor of a line bundle. Its restriction to
$E'=\C^*/\langle\tau\rangle$ is that same automorphy
line; its kernel is the discrete subgroup
$\langle\tau\rangle$. Consequently the identity
component of the restriction kernel is precisely the
compact group $p'^*\Pic^0(B')$.
Our assumed trivial restriction on $E$ makes
$\psi^*\rho(\Pic^0(Y))$ lie in that identity component.
Gauduchon degree is a continuous additive homomorphism
to $\R$, and is zero on any compact subgroup. On the
other hand this image contains
$\OO_{Z'_0}(q\psi^*A_0)$, the line of a nonzero effective
divisor, which has strictly positive Gauduchon degree.
This contradiction proves the first assertion in both
cases, and proves the stated nonconstant restriction.
Riemann--Hurwitz gives genus at most one for the
normalization of the nonconstant image of $E$, so it is
elliptic.

For the last assertion, an effective divisor of zero
real Chern class on a minimal class VII surface is,
by \cite[Remark~2.7]{DlousskyTeleman2012}, an elliptic
divisor on a Hopf surface or a positive multiple of the
rational cycle on an Enoki surface. If the source states
integral homological triviality, first multiply the
divisor by an integer killing its torsion integral
class and use Poincar\'e duality; the resulting support
and coefficient conclusion descends to the original
divisor. Formula (C.INT.2) also applies to an Enoki
minimal model. Thus the adjacent divisor on its
resolution is the full pullback of a rational cycle;
all its curve images have rational normalization.
On a Hopf or non-K\"ahler elliptic neighbouring
resolution, the adjacent divisor is the full pullback
of a nonzero divisor consisting of elliptic curves.
At least one of their strict transforms must map
into the same intersection. It cannot map to a point
or a rational-normalized curve by the result just
proved. This is the required contradiction.
\end{proof}


% C.2.INT.2 from part_C2_interfaces.tex

